\documentclass[10pt,a4paper]{amsart}
\usepackage[margin=19mm]{geometry}
\usepackage{amsmath,amssymb,mathtools}
\usepackage[scr=rsfs,cal=euler]{mathalfa}
\usepackage{xfrac}
\usepackage[unicode,hidelinks,bookmarks=false]{hyperref}
\newcommand{\ie}{i.e.\ }
\newcommand{\cf}{cf.\ }
\newcommand{\resp}{respectively\ }
\newcommand{\Subsection}{\subsection}
\providecommand{\nobreakdash}{\nobreak\hbox{-}}
\setlength{\emergencystretch}{2em}
\multlinegap0pt
\usepackage{bm}
\usepackage{bbm}
\usepackage{dsfont}
\usepackage{xspace}
\usepackage{cancel}
\usepackage{mathtools}
\usepackage{amsthm}
\makeatletter
\@ifundefined{problem}{\newtheorem{problem}{Problem}}{}
\@ifundefined{theorem}{\newtheorem{theorem}[problem]{Theorem}}{}
\makeatother
\title[Torus Queen Independence]{Torus Queen Independence}
\author{Kada Williams}
\date{Source: arXiv:2404.18237v3 (CC BY 4.0)}
\begin{document}
\maketitle

\noindent\emph{Source and licence.} Torus Queen Independence. Version arXiv:2404.18237v3 (CC BY 4.0), distributed under the Creative Commons Attribution 4.0 International licence, as recorded in the arXiv licence metadata for this exact version and shown on the abstract page https://arxiv.org/abs/2404.18237v3. Full rights notice: licenses/arxiv-2404.18237-CC-BY-4.0.txt.\par\smallskip

\noindent\textit{Editorial scope.} Counted unit: the Theorem whose source label is \texttt{525} in the article (source0.tex lines 122--124, source label \texttt{525}), delivered together with the article's own complete original proof of it (source0.tex lines 126--161, one \texttt{proof} environment of 36 lines). No mathematics is written, repaired or re-derived: statement and proof are the article's own text line for line. Required context retained uncounted: none. Every definition, display and label of those lines is retained unchanged. Omitted from this package and not redistributed: 1--121, 125--125, 162--331. Nothing inside a retained excerpt is omitted or altered: every retained body is delivered exactly as the article's lines are written, so no omission receipt of any kind accompanies this package. Citations: the retained excerpts carry no citation command at all, so no bibliography entry of the article is transcribed and no reference is redistributed. Reference adaptation: none was needed and none was made; every reference inside the delivered unit resolves inside it, so no unresolved reference or citation is produced. Printed slips kept literal: no printed word, symbol, hypothesis or number of the counted statement or of its proof was altered. Layout and class: the article's own class is \texttt{article} with packages including times, amsmath, amssymb, theoremref, inputenc, indentfirst, graphicx, tikz, xcolor, xskak, fontenc, amsthm; the wrapper is the lane's pinned \texttt{amsart} editorial wrapper at 10pt on A4, which restates the theorem-like environments the retained text needs and adds the article's own macro definitions that the retained text needs (none), with extra wrapper packages (bm, bbm, dsfont, xspace, cancel, mathtools). The delivered numbering is the unit's own. Assets: no retained span contains a figure environment, a graphics-inclusion command, a PDF-page inclusion, a table or a TikZ picture; no image file is copied into this package, the package declares no asset dependency and no image file is redistributed with it. Licence: version arXiv:2404.18237v3 of this article is distributed under the Creative Commons Attribution 4.0 International licence, as recorded in the arXiv licence metadata for this exact version and shown on the abstract page https://arxiv.org/abs/2404.18237v3. A full rights notice is delivered at licenses/arxiv-2404.18237-CC-BY-4.0.txt. Characters: every retained excerpt is pure ASCII, so the delivered engine, XeLaTeX with its default Latin Modern text font, reports no missing character.
\begin{theorem}\label{525}
Let $n$ be a positive integer that is not a multiple of $2$ or $3$, divisible by $5$, and not a multiple of $25$. For such $n$, there is no independent set of $n^{2}$ queens on $\mathbb{Z}_n^3$.
\end{theorem}

\begin{proof}
Suppose that queens are placed on
$$(x_1,y_1,z_1),\dots,(x_{n^2},y_{n^2},z_{n^2})\in \mathbb{Z}_n^3.$$
We claim that for these fields $(x,y,z)$, the pairs $(y,z)$ attain pairwise distinct values, and similarly for $(x\pm y,z)$ and $(x\pm y,x\pm z)$ with any given signs. This is the case because if two fields gave rise to the same value of a pair, then they would be apart by a multiple of $(1,0,0)$, $(1,\mp 1,0)$, or $(1,\mp 1,\mp 1)$, respectively. Hence, each of these $n^2$ pairs take on every value in $\mathbb{Z}_n^2$.

Let $k$ and $l$ be positive integer exponents, and consider the sum
$$\sum_{i=1}^{n^2} x_i^k y_i^l.$$
Since $(x_i,y_i)$ takes on every possible value precisely once, this equals
$$\sum_{(x,y)\in \mathbb{Z}_n^2}x^ky^l=\left(\sum_{x\in \mathbb{Z}_n}x^k\right)\left(\sum_{y\in \mathbb{Z}_n}y^l\right).$$

The value of a power sum $S_k=\sum_{x=0}^{n-1}x^k$ can be computed recursively, due to the well-known combinatorial identity $\sum_{x=k}^{n-1} \binom {x}{k}=\binom{n}{k+1}$, which we write as
$$\sum_{x=0}^{n-1} \frac{x(x-1)\dots(x-k+1)}{k!}=\frac{n(n-1)\dots(n-k)}{(k+1)!}.$$
Upon multiplying both sides $k!$, we can expand $x(x-1)\dots(x-k+1)$, the sum of $x^k$ and lower order terms, which, when summed over $x$, yields the sum of $S_k$ and an integer combination of $S_{k-1},\ldots,S_0$. Since $n$ is not a multiple of $2$ or $3$, it divides the integer $\frac1{k+1} n(n-1)\dots(n-k)$ for $k\le 3$. Since $n$ is a multiple of $5$, for $k=4$, this expression is $\frac{n}{5}(-1)(-2)(-3)(-4)\equiv -\frac{n}{5} \pmod n$. Thus, we deduce that $S_k\equiv 0\pmod n$ if $k\le 3$ and $S_4\equiv -\frac{n}{5}\pmod n$.

Projecting to $\mathbb{Z}_n$, these considerations motivate us to work with
$$\sum_{i=1}^{n^2} x_i^4 y_i^4=S_4\cdot S_4,$$
which does not vanish, provided that $n$ is not divisible by $25$. If we use the shorthand $\sum f(x,y,z)$ to mean $\sum_{i=1}^{n^2} f(x_i,y_i,z_i)$, then since pairs of the type $(x\pm y,z)$ and $(x\pm y,x\pm z)$ also attain every possible value,
$$\sum y^4(x\pm z)^4=S_4^2,$$
$$\sum (x\pm y)^4(x\pm z)^4=S_4^2.$$
To cancel terms, we add, subtract, and factor:
\begin{equation}
\sum \left[(x+y)^4+(x-y)^4\right]\cdot \left[(x+z)^4+(x-z)^4\right]=4\cdot S_4^2,
\end{equation}
\begin{equation}
\sum \left[(x+y)^4+(x-y)^4-2y^4\right]\cdot \left[(x+z)^4+(x-z)^4-2z^4\right]= 0.
\end{equation}
Since $(x+y)^4+(x-y)^4=2(x^4+6x^2y^2+y^4)$, (6.2) shows $\sum (x^4+6x^2y^2)(x^4+6x^2z^2)=\sum x^8+\sum x^6z^2+\sum x^6y^2+\sum x^4y^2z^2$ to vanish. Of course, we know that $\sum x^6y^2=S_6S_2$ vanishes, as well as $\sum x^6z^2=S_6S_2$ and $\sum x^8=S_8S_0$. Thus, we learn that $\sum x^4y^2z^2$ vanishes modulo $n$. Returning to (6.1),
$$\sum (x^4+6x^2y^2+y^4)(x^4+6x^2z^2+z^4)=S_4^2$$
implies that
\begin{align*}
    & \sum (x^8+6x^6z^2+6x^6y^2) \\
    + &\sum (x^4y^4+x^4z^4+y^4z^4)\\
    + & \sum (6x^2y^2z^4+6x^2y^4z^2+36x^4y^2z^2)\equiv S_4^2 \pmod n.
\end{align*}
Of these three terms, we know that the first term vanishes, and the third vanishes by virtue of $\sum x^4y^2z^2=0$ and its variants. The remaining second term is $3\cdot S_4^2$, whence we arrive at a contradiction, $3S_4^2=S_4^2$, i.e. $S_4^2=0$ in $\mathbb{Z}_n$.
\end{proof}

\end{document}
